\begin{afterword}
\summaryhead

The author recounts a conversation of a few hours with an intelligent, nominally Orthodox Jewish woman. She knew some of the evidence against her religion, for example that the characters of the Book of Esther are not in the Persian records, and she still celebrated Purim. As an atheist in high school she had decided to act as if she believed in God, and she said that over time she had come to believe.

The author concludes, ``so far as I can tell,'' that she had not. She always said ``I believe in God'' and never ``There is a God,'' she spoke only of the benefits of believing, and she agreed that someone who just wanted the truth would not invent God as a hypothesis. What she has instead, the author thinks, is ``belief in self-deception'': she believes she has deceived herself. The author supposes that this brings much the same placebo effect as real belief, which would explain why she defended ``I believe'' and showed no interest in whether God exists.

\responsehead

The idea is plausible and carefully hedged: a person can believe they have deceived themselves without having done so. The post keeps its hedges (``so far as I can tell,'' ``I think,'' ``I suppose''), and it claims only to understand something better, not to have discovered it.

The evidence is one conversation, reported from one side, and the diagnosis reads her mind from her words. The notes on ``Professing and Cheering'' raise the same concern about its panelist. Of the three pieces of evidence, the one the post puts first is the weakest. ``I believe in God'' is the ordinary form of a profession of faith. Merriam-Webster defines believing in God as regarding God's existence as a fact, and the thirteen lines of \textsc{Ani Ma'amin} in Jewish prayer each begin ``I believe with full faith.'' Her agreement that a truth-seeker would not invent God as a hypothesis is stronger evidence.

Her method, acting as if she believed until she did, is the one Pascal recommended. It is the indirect route that ``Doublethink'' left undiscussed. The post judges that in her case it failed, on the evidence of the same conversation.

The explanatory step is supposed, not shown. That believing you have deceived yourself ``ought to produce much the same placebo effect as actually believing in God'' comes with ``I suppose'' and nothing else. It also sits uneasily with ``Doublethink,'' which said you ``might even believe you were happy and self-deceived; but you would not in fact be happy and self-deceived.'' Read strictly, the two can be reconciled, but the post does not do it.

Two background claims are stated more firmly than the evidence allows. ``Most people like this'' refuse to talk to atheists, the post says, without evidence. And what ``she knows'' about Esther includes an 1890s scholarly conjecture, and leaves out Persian court records that name officials called Marduka.

In short: a hedged diagnosis of one person from one conversation, whose first piece of evidence is a standard form of words and whose explanation, a placebo from believing in one's own self-deception, is supposed rather than shown.
\end{afterword}
